\lecture{46}{Несогласованность цели}

\sourcevideo
    {https://www.youtube.com/playlist?list=PLOzRYVm0a65fhKDS8cYIC7YCuVGJ4FOJx}
    {Week 12: Lecture 46: Objective Inconsistency Problem}

В федеративном обучении клиенты могут выполнять различное число
локальных шагов. Тогда обычное усреднение изменяет не только скорость
алгоритма: уже в первом порядке оно перевзвешивает локальные функции
и может направить метод к минимуму другой целевой функции.

\section{Глобальная задача и локальные обновления}

Пусть требуется минимизировать
\[
    F(x)=\sum_{i=1}^{K}p_iF_i(x),
    \qquad
    p_i\ge0,
    \qquad
    \sum_{i=1}^{K}p_i=1.
\]
Для эмпирического риска обычно используют
\[
    p_i=\frac{n_i}{\sum_{r=1}^{K}n_r},
\]
где \(n_i\) --- число объектов у клиента \(i\).

На раунде \(t\) все выбранные клиенты начинают из общей модели
\[
    x_i^{t,0}=x^t
\]
и выполняют \(\tau_i\) локальных градиентных шагов
\[
    x_i^{t,j+1}
    =
    x_i^{t,j}-\eta\nabla F_i(x_i^{t,j}),
    \qquad
    j=0,\ldots,\tau_i-1.
\]
Обычный FedAvg формирует
\[
    x^{t+1}
    =
    \sum_{i=1}^{K}p_i x_i^{t,\tau_i}
    =
    x^t+
    \sum_{i=1}^{K}p_i\Delta_i^t,
\]
где
\[
    \Delta_i^t=x_i^{t,\tau_i}-x^t.
\]
Если значения \(\tau_i\) различаются, масштабы и направления
локальных приращений также различаются.

\section{Точный пример с двумя клиентами}

Рассмотрим
\[
    F_1(x)=(x-1)^2,
    \qquad
    F_2(x)=2(x-5)^2,
    \qquad
    p_1=p_2=\frac12.
\]
Глобальная функция равна
\[
    F(x)
    =
    \frac12(x-1)^2+(x-5)^2,
\]
поэтому
\[
    F'(x)=3x-11,
    \qquad
    x^\star=\frac{11}{3}.
\]
Локальные минимизаторы равны \(1\) и \(5\), а их обычное среднее
равно \(3\). Следовательно, усреднение локальных решений не
эквивалентно минимизации средней функции.

Для первого клиента
\[
    x_1^{t,j+1}-1
    =
    (1-2\eta)(x_1^{t,j}-1),
\]
откуда
\[
    x_1^{t,\tau_1}
    =
    1+(1-2\eta)^{\tau_1}(x^t-1).
\]
Для второго клиента аналогично
\[
    x_2^{t,\tau_2}
    =
    5+(1-4\eta)^{\tau_2}(x^t-5).
\]
Локальные итерации устойчивы при \(0<\eta<1\) для первой функции и
при \(0<\eta<1/2\) для второй, поэтому далее достаточно считать
\[
    0<\eta<\frac12.
\]

\section{Серверная динамика и её предел}

Обозначим
\[
    a_1=(1-2\eta)^{\tau_1},
    \qquad
    a_2=(1-4\eta)^{\tau_2}.
\]
Тогда серверное усреднение задаёт аффинную рекурсию
\[
    x^{t+1}
    =
    3+
    \frac{a_1}{2}(x^t-1)
    +
    \frac{a_2}{2}(x^t-5).
\]
Коэффициент при \(x^t\) равен
\[
    q=\frac{a_1+a_2}{2}.
\]
При \(0<\eta<1/2\) выполнено
\[
    |a_1|<1,
    \qquad
    |a_2|<1,
\]
а значит, \(|q|<1\). Поэтому последовательность сходится к
единственной неподвижной точке
\[
    \widetilde x
    =
    \frac{6-a_1-5a_2}{2-a_1-a_2}.
\]
В общем случае \(\widetilde x\ne x^\star\), причём смещение зависит
не только от функций, но и от \(\eta\), \(\tau_1\) и \(\tau_2\).

При \(\eta\to0\) и фиксированных \(\tau_1,\tau_2\)
\[
    a_1=1-2\eta\tau_1+O(\eta^2),
    \qquad
    a_2=1-4\eta\tau_2+O(\eta^2),
\]
откуда
\[
    \widetilde x
    \longrightarrow
    \frac{\tau_1+10\tau_2}{\tau_1+2\tau_2}.
\]
Это значение совпадает с \(11/3\) при
\(\tau_1=\tau_2\), но обычно отличается от глобального минимума.

\section{Эффективная целевая функция}

Пусть градиенты \(F_i\) липшицевы, шаг \(\eta\) мал, а числа
локальных итераций фиксированы. Тогда
\[
    x_i^{t,\tau_i}
    =
    x^t
    -
    \eta\tau_i\nabla F_i(x^t)
    +
    O(\eta^2\tau_i^2).
\]
После усреднения
\[
\begin{aligned}
    x^{t+1}
    ={}&
    x^t
    -
    \eta
    \sum_{i=1}^{K}
    p_i\tau_i\nabla F_i(x^t)
    \\
    &+
    O\left(
        \eta^2
        \sum_{i=1}^{K}p_i\tau_i^2
    \right).
\end{aligned}
\]
Положим
\[
    s=\sum_{r=1}^{K}p_r\tau_r,
    \qquad
    \widetilde p_i=\frac{p_i\tau_i}{s}.
\]
Тогда главный член является градиентным шагом размера \(\eta s\) по
функции
\[
    \widetilde F(x)
    =
    \sum_{i=1}^{K}\widetilde p_iF_i(x).
\]
Отличие \(\widetilde F\) от требуемой функции \(F\) называется
несогласованностью целевой функции. При весах \(p_i=n_i/n\)
эффективный вес клиента пропорционален \(n_i\tau_i\), а не только
числу его объектов.

Эта интерпретация является локальной по шагу \(\eta\). При конечном
шаге и нелинейных функциях градиенты вычисляются в промежуточных
точках \(x_i^{t,j}\), поэтому обновление FedAvg не обязано быть точным
градиентным шагом по одной фиксированной функции. Если \(\tau_i\)
меняются от раунда к раунду, меняются и эффективные веса.

\section{Когда смещение исчезает}

Если все клиенты выполняют одинаковое число шагов,
\[
    \tau_i=\tau,
\]
то
\[
    \widetilde p_i
    =
    \frac{p_i\tau}{\sum_r p_r\tau}
    =p_i.
\]
Следовательно, в первом порядке эффективная и исходная целевые
функции совпадают. Однако равное число шагов не выравнивает время
работы: синхронный сервер по-прежнему ждёт самых медленных клиентов.

Если все локальные функции имеют одно и то же множество минимизаторов,
перевзвешивание может не изменить предельное решение, хотя изменяет
траекторию и скорость. В частности, при
\[
    F_1=\cdots=F_K
\]
любые положительные веса задают одну и ту же функцию с точностью до
множителя. Наиболее заметное смещение возникает при сочетании
неодинаковых \(\tau_i\) и статистически неоднородных локальных задач.

\section{Коррекция локальной работы}

В простейшем случае постоянного локального шага можно изменить веса
серверной агрегации. Если
\[
    q_i
    =
    \frac{p_i/\tau_i}{\sum_{r=1}^{K}p_r/\tau_r},
\]
то
\[
    q_i\tau_i
    =
    \frac{p_i}{\sum_r p_r/\tau_r}.
\]
Поэтому главный член обновления имеет требуемое направление
\[
    -\sum_{i=1}^{K}p_i\nabla F_i(x^t)
\]
с точностью до общего множителя, который можно включить в серверный
шаг.

Одного деления на \(\tau_i\) недостаточно, если клиент использует
меняющиеся шаги, импульс или другой локальный оптимизатор. Тогда
приращение представляется как взвешенная сумма локальных градиентов,
а нормирующий коэффициент должен учитывать всю локальную рекурсию.
Эта идея приводит к нормализованной агрегации FedNova, изучаемой в
следующей лекции.
